% Part: first-order-logic
% Chapter: model-theory
% Section: theory-of-m

\documentclass[../../../include/open-logic-section]{subfiles}

\begin{document}

\olfileid{mod}{bas}{thm}

\section{د يوه \printtoken{S}{structure} تيوري}

هر !!{structure}~$\Struct{M}$ ځينې !!{sentence}s رښتونې او ځينې
دروغې کوي۔ د هغو ټولو !!{sentence}s سټ چې دا جوړښت ئې رښتونې کوي، د
هغه \emph{تيوري} بلل کېږي۔ دا سټ په رښتيا يوه تيوري ده، ځکه د هغې هر
معنايي استلزام بايد د هغې په ټولو مدلونو، او له دې ډلې
په~$\Struct{M}$ کښې، رښتونے وي۔

\begin{defn}
  يو !!a{structure}~$\Struct M$ راکړل شي۔ د~$\Struct{M}$
  \emph{تيوري} د هغو !!{sentence}s سټ~$\Theory{M}$ دے چې
  په~$\Struct{M}$ کښې رښتونې وي، يعنې
  $\Theory{M} = \Setabs{!A}{\Sat{M}{!A}}$۔
\end{defn}

د ``تيورۍ'' اصطلاح په غيرصوري ډول د هغو !!{sentence}s سټونو دپاره
هم کاروو چې يو ټاکلے مطلوب تفسير لري، که د استنتاج لاندې تړلي وي او
که نۀ۔

\begin{prop}
د هر~$\Struct{M}$ دپاره $\Theory{M}$ بشپړه ده۔
\end{prop}

\begin{proof}
د هرې !!{sentence}~$!A$ دپاره يا $\Sat{M}{!A}$ وي يا
$\Sat{M}{\lnot !A}$؛ نو يا $!A \in \Theory{M}$ وي يا
$\lnot !A \in \Theory{M}$۔
\end{proof}

\begin{prop}\ollabel{prop:equiv}
  که د هر~$!A \in \Theory{M}$ دپاره $\Struct{N} \models !A$ وي، نو
  $\Struct{M} \elemequiv \Struct{N}$۔
\end{prop}

\begin{proof}
دا چې د ټولو $!A \in \Theory{M}$ دپاره $\Sat{N}{!A}$ دے،
$\Theory{M} \subseteq \Theory{N}$۔ که $\Sat{N}{!A}$ وي، نو
$\Sat/{N}{\lnot !A}$؛ ځکه نو $\lnot !A \notin \Theory{M}$۔ دا چې
$\Theory{M}$ بشپړه ده، $!A \in \Theory{M}$۔ نو
$\Theory{N} \subseteq \Theory{M}$، او له دې
$\Struct{M} \elemequiv \Struct{N}$ لرو۔
\end{proof}

\begin{rem}\ollabel{remark:R}
  $\Struct{R} = \langle\Real, <\rangle$ وګورئ: دا هغه !!{structure}
  دے چې د تعريف ساحه ئې د حقيقي عددونو سټ~$\Real$ او !!{language} ئې
  يوازې يو دوه ځاييز !!{predicate} لري چې پر حقيقي عددونو د~$<$
  اړيکې په توګه تفسير شوے۔ څرګنده ده چې $\Struct{R}$
  !!{nonenumerable} دے؛ خو دا چې $\Theory{R}$ څرګنده سازګاره ده، د
  لوېنهايم--سکولم د قضيې له مخې !!a{enumerable} مدل لري، مثلاً
  $\Struct{S}$، او د~\olref{prop:equiv} له مخې $\Struct{R}
  \equiv \Struct{S}$۔ سربېره پر دې، دا چې $\Struct{R}$ او
  $\Struct{S}$ آيزومورف نۀ دي، ښيي چې د~\olref[iso]{thm:isom} معکوس
  په عام ډول نۀ صدق کوي۔
\end{rem}

\end{document}
